\newcommand{\slidemode}{1}
% Modified by LianTze Lim to work with fontspec/xelatex
    % Choose for notes or slides
    \if \slidemode 1
        \documentclass{beamer}
        \else
        \documentclass{article}
        \usepackage{beamerarticle}
        \usepackage[a4paper,margin=2cm]{geometry}
        \usepackage{multicol}
    \fi 
\usepackage{amsfonts,amsmath,amssymb,amsthm,amscd}
\usepackage[style=authoryear]{biblatex}
\renewcommand*{\nameyeardelim}{\addcomma\addspace}
\AtBeginBibliography{\tiny}

\usetheme{DarkConsole}
\setbeamertemplate{theorems}[normal font]

% Packages
\usepackage{hyperref}
\usepackage{graphicx}
\usepackage{multimedia} % provides \movie
\usepackage{tikz,pgfplots}
\usetikzlibrary{arrows.meta,calc}
\pgfplotsset{compat = 1.18}
\usepgfplotslibrary{fillbetween}
\usepackage{physics}
\usepackage{siunitx}
\usepackage{polynom}
\usepackage{biblatex}
\usepackage[outline]{contour} % glow around text
\usetikzlibrary{angles,quotes,intersections,fillbetween} % for pic
\newcommand{\myscale}{0.7}
\contourlength{1.3pt}

\tikzset{>=latex} % for LaTeX arrow head
\usepackage{xcolor}
\tikzstyle{vector}=[->,very thick,xcol]
\tikzstyle{mydashed}=[dash pattern=on 2pt off 2pt]
\def\tick#1#2{\draw[thick] (#1) ++ (#2:0.1) --++ (#2-180:0.2)} 

% helper: open-dot macro (hollow circle)
  \newcommand{\opendot}[2]{%
    \filldraw[fill=white,draw=gray,thick] (#1,0) circle[radius=2.2pt];
    \node[below=4pt] at (#1,0) {#2};
  }
  \newcommand{\closedot}[3][magenta]{%
    \filldraw[fill=#1,draw=#1,thick] (#2,0) circle[radius=2.2pt];
    \node[below=4pt] at (#2,0) {#3};
  }
\usepackage[brazilian]{babel}
\usepackage{diagbox}
\usepackage{booktabs}

% collor-blind friendly pallete
\definecolor{figcolor1}{RGB}{0,70,70}
\definecolor{figcolor2}{RGB}{116,233,233}
\definecolor{figcolor3}{RGB}{146,0,0}
\definecolor{color1}{RGB}{0,70,70}
\definecolor{color2}{RGB}{116,233,233}
\definecolor{color3}{RGB}{146,0,0}


% A counter, since TikZ is not clever enough (yet) to handle
% arbitrary angle systems.
\newcount\mycount
\usepackage[percent]{overpic}
\usepackage{xcolor}
\usepackage{adjustbox}
\usepackage{bbold}
\newcommand{\rest}{\upharpoonright}

% symbols
\newcommand{\F}{\mathcal{F}}
\newcommand{\C}{\mathcal{C}}
\newcommand{\Open}{\mathcal{O}}
\newcommand{\A}{\mathcal{A}}
\newcommand{\B}{\mathcal{B}}
\newcommand{\G}{\mathcal{G}}
\newcommand{\M}{\mathcal{M}}
\newcommand{\N}{\mathbb{N}}
\newcommand{\Q}{\mathbb{Q}}
\newcommand{\meta}[1]{\overset{\text{meta}}{#1}}
\newcommand{\val}[1]{\llbracket #1 \rrbracket}
\newcommand{\valin}[3]{
	\sup_{#3 \in dom(#2)}\val{#3 = #1}#2(#3)}
\newcommand{\valsub}[3]{
	\inf_{#3 \in dom(#1)}(#1(#3)\implies \val{#3 \in #2})}
\newcommand{\valeq}[2]{\val{#1 \subset #2}\val{#2 \subset #1}}
% epimorphism arrow
\newcommand{\rightarrowdbl}{\rightarrow\mathrel{\mkern-14mu}\rightarrow}
\newcommand{\epimorphism}[2][]{%
	\xrightarrow[#1]{#2}\mathrel{\mkern-14mu}\rightarrow
}
\newcommand{\catname}[1]{\textnormal{\textsc{#1}}}

\newcommand{\R}{\mathbb{R}}
\newcommand{\D}{\mathcal{D}}
\newcommand{\Part}{\mathbb{P}}
\newcommand{\sequence}[1]{\langle #1 \rangle}
\DeclareMathOperator{\inter}{int}
\newcommand{\interior}[2]{\inter_{#2} (#1)}
\DeclareMathOperator{\clos}{cl}
\DeclareMathOperator{\diam}{diam}
\newcommand{\closure}[2]{\clos_{#2} (#1)}
\newcommand{\st}{\,:\,}
\newcommand{\dem}{\emph{dem.:}\\}
\renewcommand{\d}{\mathrm{d}}

% Theorem environments
\numberwithin{equation}{section}
% Theorem environments
\newtheorem{teor}{Teorema}[section]
\newtheorem{lema}{Lema}[section]
\newtheorem{prop}{Proposição}[section]
\newtheorem{properties}{Propriedades}[section]
\newtheorem{cor}{Corolário}[section]  
\newtheorem{defin}{Definição}[section]
\newtheorem{ex}{Exemplo}[section]
\newtheorem*{theorem*}{Teorema}
\newtheorem*{proposition*}{Proposição}
\newtheorem*{corollary*}{Corolário}
\newtheorem*{lemma*}{Lema}
\theoremstyle{remark}
\newtheorem{rmk}{Observação}[section]
\newtheorem*{question}{Problema}
\newtheorem{exerc}{Exercício}

\title{\texttt{Álgebras de Boole de grafos infinitos (e espaços de emaranhados!)}}
\author{Violeta Mar\\ Orientado por Leandro Aurichi \\ Lab Top, ICMC}
\date{}
\addbibresource{bibliography.bib}

\begin{document}

\begin{frame}{}
\maketitle
\end{frame}

\begin{frame}{}
\frametitle{Introduções}
    \begin{itemize}
        \item Vou apresentar objetos, ideias e resultados do artigo que escrevi com o Gustavo Boska ano passado. \pause
        \item Trabalhamos em Combinatória Infinita, mais especificamente em grafos infinitos. \pause
        \item Vou tentar fazer uma apresentação para quem não é da área, então vai ter bastante definição. \pause
        \item No final, vou dar uma ideia do que estou trabalhando este ano. 
    \end{itemize}
\end{frame}

\begin{frame}{}
\frametitle{Extremidades}
    \begin{itemize}
        \item Um grafo consiste de um conjunto de vértices $V$ e um conjunto de arestas $E$ entre vértices (uma relação simétrica). \pause
        \item Um grafo é infinito quando $V$ é infinito. \pause
        \item Um caminho é uma sequência de vértices distintos $(v_1, v_2, \dots)$ onde $v_i$ é adjacente a $v_{i+1}$. \pause
        \item Um raio é um caminho infinito. \pause
        \item Dois raios são ditos equivalentes se existem infinitos caminhos disjuntos ligando eles entre si. \pause
        \item Uma extremidade é a classe de equivalência de caminhos infinitos sob essa equivalência.
        \item O conjunto das extremidades é denotado como $\Omega(G)$ e ele admite uma topologia natural.
    \end{itemize}
\end{frame}
\begin{frame}{}
\frametitle{Direções (= extremidades!)}
    \begin{itemize}
        \item Uma direção $\rho$ em um grafo $G$ é uma função que, dado um conjunto finito de vértices $F$ te retorna um conjunto de vértices $\rho(F)$ que corresponde a uma componente conexa do grafo $G \setminus F$. \pause
        \item Uma direção deve satisfazer:  $F \subset F'$ então $\rho(F') \subset \rho(F)$ \pause
        
        \begin{theorem}[Diestel e Kuhn, 2003]
        Toda extremidade define uma direção e toda direção vem de uma extremidade dessa forma.
         \end{theorem}
         \pause
         \newline
        \item Dado um raio $r$, \ direção $\rho_r$ é a que seleciona, dado um conjunto finito de vértices $F$, a componente $\rho(F)$ que tem infinitos vértices de $r$.\pause
        \item A topologia em $\Omega(G)$ é a gerada pela base de conjuntos da forma $\{r \mid \rho_r(F) = \rho(F)\}$ para um dado conjunto finito $F$ de vértices e uma direção $\rho$.

        \end{itemize}
\end{frame}

\begin{frame}{}
\frametitle{Emaranhados}
    \begin{itemize}
        \item Dado um subconjunto de vértices $A \subset V$, o grafo $G[A]$ é o subgrafo de $G$ contendo todos os vértices em $A$ com todas as arestas entre eles.\pause
        \item Uma separação $\{A,B\}$ é composta de dois subconjuntos de vértices $A,B$ tal que $G = G[A] \cup G[B]$ \pause
        \item A ordem de uma separação é a cardinalidade $|A \cap B|$ \pause
        \item Uma orientação $\tau$ é uma escolha de orientação $(A,B)$ para cada separação $\{A,B\}$ de ordem finita de um dado grafo. \pause
        \item Dadas $n$ separações ordenadas $(A_1,B_1),\dots, (A_n, B_n)$, dizemos que elas formam uma $n$-contradição se $G[A_1] \cup \dots G_[A_n] = G$. \pause
        \item Um emaranhado é uma orientação sem $1,2,3$-contradições.\pause
        \begin{theorem}[Diestel, 2017]
        Toda extremidade define um emaranhado. Mais: os emaranhados possuem uma topologia natural onde o espaço de extremidades se mergulha. O espaço de emaranhados é compacto, Hausdorff e zero-dimensional. (de Stone!)
         \end{theorem}

        \end{itemize}
\end{frame}

\begin{frame}{}
\frametitle{Aresta-Extremidades e Aresta-Direções}
    \begin{itemize}
        \item Dois raios são ditos aresta-equivalentes se existem infinitos caminhos disjunto-de-arestas ligando eles entre si. \pause
        \item Uma aresta-extremidade é a classe de equivalência de caminhos infinitos sob essa aresta-equivalência.\pause
        \item Uma aresta-direção $\rho$ em um grafo $G$ é uma função que, dado um conjunto finito de arestas $F_E$ te retorna um conjunto de vértices $\rho(F_E)$ que corresponde a uma componente conexa do grafo $G \setminus F_E$. \pause
        \item Uma aresta-direção deve satisfazer:  $F_E \subset F_{E}'$ então $\rho(F_{E}') \subset \rho(F_E)$ \pause
        \item Definimos topologias nas aresta-extremidades e nas aresta-direções analogamente ao caso dos vértices.
        \begin{theorem}[Boska, Duzi, Magalhães - 2025]
            O espaço de aresta-extremidades se mergulha naturalmente no espaço de aresta-direções. O espaço de aresta-direções é compacto, Hausdorff e zero-dimensional. (de Stone!)
        \end{theorem}
        \end{itemize}
\end{frame}

\begin{frame}{}
\frametitle{Aresta-Emaranhados?}
    \begin{itemize}
        \item Se definirmos um análogo dos emaranhados para arestas, vamos apenas obter o espaço de aresta-direções novamente. 
    \end{itemize}
\end{frame}

\begin{frame}{}
\frametitle{A situação}
%desenhar grafo G -> omega(G) contido dentro de T(G) de Stone e omega_E(G) contido dentro de T_E(G) de Stone
\end{frame}

\begin{frame}{}
\frametitle{Dualidade de Stone}
    \begin{itemize}
        \item Álgebras de Boole são estruturas algébricas $(\mathcal{B}, \vee, \wedge, \lnot)$ que modelam as operações de união (ou), interseção (e), complementar (negação). \pause
        \item Um subconjunto de $\mathcal{P}(X)$ fechado para união, interseção e complementar é uma álgebra de Boole \pause (na verdade, esses são todos os exemplos...) \pause
        \item Espaços de Stone são espaços topológicos compactos, Hausdorff e zero-dimensionais. \pause
        \begin{theorem}[Dualidade de Stone]
            A categoria de álgebras de Boole é equivalente a categoria de espaços de Stone.
        \end{theorem}
    \end{itemize}
\end{frame}


\begin{frame}{}
\frametitle{Álgebra -> Topologia}
    \begin{itemize}
        \item Um ultrafiltro em uma álgebra de Boole $\mathcal{B}$ é um homomorfismo $u: \mathcal{B} \rightarrow \{0,1\}$ \pause
        \item O conjunto $\text{Ult}(\mathcal{B})$ admite uma topologia natural, dada pela base de conjuntos da forma $\{u \mid u(b) = 1\}$ para cada $b \in \mathcal{B}$ \pause
        \item $\text{Ult}(\mathcal{B})$ é um espaço de Stone \pause
        \item $\text{Ult}$ é funtorial.
    \end{itemize}
\end{frame}

\begin{frame}{}
\frametitle{Topologia -> Álgebra}
    \begin{itemize}
        \item Dado um espaço de Stone $X$, o conjunto $\text{Clp}(X) = \{U \subset X \mid U \text{ é aberto e fechado}\}$ é fechado para união, interseção e complementar - uma álgebra de Boole \pause
        \item $\text{Clp}$ é funtorial. \pause
        \item $\text{Ult}$ e $\text{Clp}$ se invertem: $\text{Ult}(\text{Clp}(X)) \cong X$  e $\text{Clp}(\text{Ult}(\mathcal{B})) \cong \mathcal{B}$
    \end{itemize} 
\end{frame}

\begin{frame}{}
\frametitle{A álgebra de separações e a álgebra de cortes}
    \begin{itemize}
        \item O conjunto $\mathfrak{s}(G) = \{A \subset E(G) \mid V(A) \cap V(E(G) \setminus A) \text{ é finito}\}$ é fechado para união, interseção e complementar. O chamamos de álgebra de separações de $G$. \pause
        \item O conjunto $\mathfrak{c}(G) = \{A \subset V(G) \mid \text{existem finitas arestas entre } A \text{ e } G \setminus A\}$ é fechado para união, interseção e complementar. O chamamos de álgebra de cortes de $G$. \pause
    \begin{theorem}[Violeta e Gustavo, 2026]
        O espaço de ultrafiltros de um quociente da álgebra de separações é homeomorfo ao espaço de emaranhados.\\
        O espaço de ultrafiltros de um quociente da álgebra de cortes é homeomorfo ao espaço de aresta-emaranhados (aresta-direções).
    \end{theorem}
    \end{itemize} 
\end{frame}

\begin{frame}[fragile]
\frametitle{Essas associações são funtoriais?}
% https://tikzcd.yichuanshen.de/#N4Igdg9gJgpgziAXAbVABwnAlgFyxMJZABgBpiBdUkANwEMAbAVxiRAB12cYAPHAIwBmwAOIAnOmgAWAXxAzS6TLnyEUAFnJVajFm07c+Q4ACEIEBnIVLseAkQBMpAIzb6zVog5deA4QBUINCttGCgAc3giUEExCABbJDIQHAgkZ2p3PS8DX2AAQ4ZI-gkAAlhSuBg0OgkAc4BX+CtFEFiEpOpUpCcdD30fPmB4GrqIcphSmHjaujApOihMOWoGOn4YBgAFZTs1EDEscKkceVb2xMQMlLTEXqzPb0McYbhR8YqmBhwJQSxvuJwEIyIA
\begin{center}
\begin{tikzcd}
\textbf{Graph} \arrow[rrrr, "\text{álgebra de separações}"] \arrow[rrd, "\text{espaço de emaranhados}"'] &  &              &  & \textbf{Bool} \arrow[lld, "\text{espaço de ultrafiltros}"] \\
                                                                                                         &  & \textbf{Top} &  &                                                           
\end{tikzcd}
\end{center} \pause 

\begin{itemize}



\item Será que um morfismo de grafos se torna um homomorfismo entre suas álgebras de separações? E uma função contínua entre seu espaço de emaranhados? \pause


\item Invariantes funtoriais são famosamente úteis. Poderiamos determinar obstruções para construções combinatoriais. \pause

\item Conseguimos funtorializar, encontrando uma classe restrita de morfismos de grafos.
\end{itemize}
\end{frame}

\begin{frame}{}
\frametitle{Emaranhados e ultrafiltros}

\begin{itemize}
    \item Emaranhados tem muito a cara de ultrafiltros: uma escolha de orientações consistentes. \pause
    \item Essas escolhas não são feitas numa álgebra de Boole, mas em um `conjunto de separações' \pause
    \item Como esse `conjunto de separações' não é fechado para interseção, os emaranhados meio que são `quase homomorfismos' \pause
    \item A teoria abstrata de emaranhados em `conjuntos de separações' arbitrários teve bastante sucesso na combinatória finita, até com aplicações em ciência de dados \pause
    \item Estou traduzindo e adaptando essa teoria para o mundo de teoria de conjuntos e lógica, onde emaranhados se tornam, em algum sentido, uma generalização de ultrafiltros que permitem o desenvolvimento de um tipo de `lógica paraconsistente'

\end{itemize}

\end{frame}



\end{document}
